\clearpage\section{MOS transistor schematic}
\subsection{CMOS inverter with independent bulk terminals}
\begin{center}
\begin{circuitikz}[font=\small, line width=.65pt]
\node[pmos,bulk] (P) at (2,4) {};
\node[nmos,bulk] (N) at (2,1) {};
\draw (P.S) -- ++(0,.65) node[above] {$vdd$};
\draw (N.S) -- ++(0,-.5) node[ground] {};
\draw (P.D) -- (N.D);
\draw (P.D |- 0,2.5) node[circ] {} -- (4.8,2.5) node[right] {$y$};
\draw (N.G) -- (.1,1) -- (.1,2.5) node[circ] {} -- (-.8,2.5) node[left] {$a$};
\draw (P.G) -- (.1,4) -- (.1,2.5);
\draw (P.bulk) -- ++(.8,0) node[right] {$bp$};
\draw (N.bulk) -- ++(.8,0) node[right] {$bn$};
\node[anchor=east] at (1.6,4.6) {\texttt{MP}};
\node[anchor=east] at (1.6,.4) {\texttt{MN}};

\end{circuitikz}
\par\small INV\_BODY in \nolinkurl{01_ring_oscillators.sp}; bp and bn are separate from the source supply rails.
\end{center}
$W_p/W_n/L=1.8/0.6/0.25\,\mu$m. With $vdd=VS=1$ V, $bp=VS+PB$ and $bn=-NB$. Cases use $(PB,NB)=(0,0),(0.2,0),(0,0.2),(0.2,0.2)$ V. All stages use the same sizes and bulk biases, with no added output load capacitor.

\subsection{Complete three-stage ring transistor wiring}
\begin{center}
\begin{circuitikz}[font=\small,line width=.65pt,scale=.9,transform shape]
\foreach \xx/\i in {0/0,4/1,8/2}{
 \begin{scope}[shift={(\xx,0)}]
 \node[pmos,nobulk] (P\i) at (1.6,4) {};
 \node[nmos,nobulk] (N\i) at (1.6,1) {};
 \draw (P\i.S) -- ++(0,.55) node[above] {$vdd$};
 \draw (N\i.S) -- ++(0,-.4) node[ground] {};
 \draw (P\i.D) -- (N\i.D);
 \draw (P\i.G) -- (0,4) -- (0,1) -- (N\i.G);
 \coordinate (g\i) at (0,2.5);
 \draw (P\i.D |- 0,2.5) node[circ] {} -- (2.8,2.5) coordinate (o\i);
 \node[above] at (2.8,2.5) {$r3\_\i$};
 \node at (.3,-.65) {$Xr3\_\i$};
 \end{scope}
}
\draw (o0) -- (g1) node[circ] {};
\draw (o1) -- (g2) node[circ] {};
\draw (o2) -- (11.8,2.5) -- (11.8,-1.3) -- (-.7,-1.3) -- (-.7,2.5) -- (g0) node[circ] {};
\end{circuitikz}
\par\small For every stage, PMOS B=bp and NMOS B=bn. Outputs r3\_0, r3\_1 and r3\_2 are cascaded, with r3\_2 returned to the first input.
\end{center}
\clearpage

\subsection{Five/seven-stage connectivity and threshold fixtures}
\begin{center}
\begin{circuitikz}[font=\scriptsize,line width=.6pt,scale=.9,transform shape]
\foreach \N/\yy in {5/3,7/0}{
 \pgfmathtruncatemacro{\last}{\N-1}
 \foreach \i in {0,...,\last}{
  \node[draw,minimum width=1.25cm,minimum height=.7cm] (q\N\i) at (1.8*\i,\yy) {$Xr\N\_\i$};
  \node[above] at (1.8*\i+.75,\yy+.35) {$r\N\_\i$};
  \ifnum\i>0\pgfmathtruncatemacro{\prev}{\i-1}\draw[->] (q\N\prev.east) -- (q\N\i.west);\fi
 }
 \draw[->] (q\N\last.east) -- ++(.35,0) -- ++(0,-1) -| ([xshift=-.35cm]q\N0.west) -- (q\N0.west);
}
\end{circuitikz}
\par\small Each block is the two-MOS INV\_BODY shown above, directed left to right; the last output returns to stage 0.
\end{center}
The three rings coexist without signal connections between them; only vdd, 0, bp and bn are shared. Each $rN\_0$ has an initial-condition assignment of 0.05 V with UIC; this is not an added physical voltage source.

\begin{center}
\begin{circuitikz}[font=\small,line width=.65pt]
\node[nmos,bulk] (N) at (2,2) {};
\draw (N.D) -- ++(0,.7) node[above] {$vdd=1\,\mathrm V$};
\draw (N.G) -- ++(-.8,0) node[left] {$vdd$};
\draw (N.bulk) -- ++(.6,0) node[right] {$0$};
\draw (N.S) -- ++(0,-.4) coordinate (sx) node[right] {$sx$};
\draw (sx) to[V,l_={\texttt{VSRC}}] ++(0,-2) node[ground] {};
\node[pmos,bulk] (P) at (8,2) {};
\draw (P.D) -- ++(0,-.5) node[ground] {};
\draw (P.G) -- ++(-.8,0) node[left] {$0$};
\draw (P.bulk) -- ++(.6,0) node[right] {$vdd$};
\draw (P.S) -- ++(0,.65) node[above] {$sp=1-V(sx)$};
\node at (1.1,2.75) {\texttt{MN}};
\node at (7.1,2.75) {\texttt{MP}};
\end{circuitikz}
\par\small MN and MP in \nolinkurl{02_body_threshold.sp}. Behavioral source EPN sets the PMOS source to $1-V(sx)$; the NMOS/PMOS bulks stay at 0/vdd.
\end{center}
VSRC sweeps from 0 to 1 V in 0.01 V increments. MN has D=G=vdd, S=sx, B=0; MP has D=G=0, S=sp, B=vdd. Sizes are the same as on the previous page; VTH is read from each MOS.

